\section{Sealing Transactions with \leanseal}\label{sec:mempool}

This section treats the puzzle as a black box and describes how it backs an encrypted mempool. The sender runs $\Build$ to get a puzzle $Z$, which contains a commitment $C$, and a key $K$. Anyone can check a claimed key against $C$ with one hash. Anyone can also recover $K$ from $Z$ alone with $N$ sequential hashes, which takes at least $T = N\ell$ seconds on the fastest hardware. \Cref{sec:construction,sec:security} construct the puzzle and prove this.

\subsection{The flow}\label{sec:flow}

\leanseal is used as a timed commitment~\cite{BN00}; \cref{fig:commit-reveal} shows the flow.
\begin{enumerate}[leftmargin=1.5em]
  \item \textbf{Submit.} The sender takes a puzzle from its stock and encrypts its signed transaction under $K$. It sends the sealed transaction together with a plaintext envelope: the account that pays, a nonce, a gas limit, a fee and a deposit (\cref{sec:admission}).
  \item \textbf{Order.} The builder orders sealed transactions without reading them. The ordering commits to the hash of each whole sealed transaction, not just to $C$: once $K$ is known, anyone could seal a different transaction under the same puzzle.
  \item \textbf{Reveal.} Once the ordering is fixed, the sender publishes $K$. Nodes check it against $C$, decrypt in about 12\,ms and execute the transaction in its committed position. Nobody solves anything.
  \item \textbf{Fallback.} If the sender stays silent past the reveal deadline, its deposit is forfeited, and anyone can solve the puzzle and publish $K$. A well-formed puzzle therefore opens without its sender, but its transaction can only run late or not at all.
\end{enumerate}

% Commit-reveal with a time-lock fallback (Figure 3 of the web version).
\begin{figure}[t]
  \centering
  \begin{tikzpicture}[
      x=1cm, y=1cm, font=\small,
      box/.style={draw, rounded corners=3pt, minimum height=1.1cm, align=center, fill=black!3},
      sub/.style={font=\footnotesize, text=black!60},
      lab/.style={font=\footnotesize, align=center},
      arr/.style={-{Stealth[length=5pt]}, line width=0.6pt},
    ]
    \node[box, minimum width=2.8cm] (wallet) at (1.4,0)
      {\textbf{Sender's wallet}\\[-1pt]{\footnotesize\color{black!60}precomputed puzzle}};
    \node[box, minimum width=2.6cm] (builder) at (7.3,0)
      {\textbf{Builder}\\[-1pt]{\footnotesize\color{black!60}orders it blind}};
    \node[box, minimum width=2.6cm] (nodes) at (11.4,0)
      {\textbf{Nodes}\\[-1pt]{\footnotesize\color{black!60}check $\mathsf{H}(K) = C$,}\\[-2pt]{\footnotesize\color{black!60}execute}};
    \node[box, dashed, fill=none, minimum width=3.6cm] (solver) at (7.3,-2.5)
      {\textbf{Any solver}\\[-1pt]{\footnotesize\color{black!60}runs $n \cdot m$ steps, publishes $K$}};
    % Main path.
    \draw[arr] (wallet) -- node[lab, above] {$\mathrm{Enc}_K(\mathit{tx})$, puzzle,\\$C = \mathsf{H}(K)$}
                          node[lab, below, text=black!60] {fee paid\\in the clear} (builder);
    \draw[arr] (builder) -- node[lab, above] {order\\fixed} (nodes);
    % Normal path: the user reveals K.
    \draw[arr, accent, line width=0.9pt]
      (wallet.north) .. controls (2.0,2.1) and (10.6,2.1) .. (nodes.north)
      node[pos=0.5, above, font=\footnotesize\bfseries, text=accent] {normal path: sender reveals $K$ right after ordering};
    % Fallback: the user stays silent.
    \draw[arr, dashed] (builder.south) -- node[lab, right, text=black!60] {key withheld} (solver.north);
    \draw[arr, dashed] (solver.east) .. controls (10.6,-2.5) and (11.4,-1.8) .. (nodes.south)
      node[pos=0.3, below right, lab] {after sequential recovery};
    % Side note.
    \node[sub, align=left, anchor=west] at (12.9,0) {No shared solve:\\one solve per\\withheld key};
  \end{tikzpicture}
  \caption{Commit--reveal with a time-lock fallback. When senders reveal their keys, no puzzle is solved. For a well-formed puzzle, withholding the key cannot prevent disclosure, but it can still delay or cancel execution.}
  \label{fig:commit-reveal}
\end{figure}


A ciphertext bound to a commitment to its key is the shape that recent Ethereum designs already use. In LUCID (EIP-8184)~\cite{LUCID}, a sealed transaction's ticket goes on chain and carries a commitment $\mathsf{keccak256}(k_{\mathrm{dem}})$ to the key and a hash of the ciphertext envelope, which itself travels off chain. The key is released by a key publisher of the sender's choice, possibly the sender itself, and the publication fee is paid only when decryption succeeds. In the hybrid design of Wagner and Ma~\cite{WM25}, the sender puts $\mathrm{Enc}_k(\mathit{tx})$ and $\mathsf{H}(k)$ on chain and gives a committee a threshold encryption of $k$, so that either the sender or the committee can open. \leanseal fits both as a fallback that needs no third party: the sender is its own key publisher, and the puzzle replaces the committee as the party that discloses the key when the sender does not. Inside LUCID only the puzzle is used. It travels in the envelope's opaque header with $K$ as $k_{\mathrm{dem}}$; its SHA-256 commitment $C$ is separate from the ticket's keccak commitment, and the cipher stays ChaCha20-Poly1305. The puzzle provides eventual disclosure only: executing a transaction whose key is disclosed minutes late needs an additional protocol rule, which LUCID does not have.

The reveal should be due as soon as the ordering is fixed. In LUCID, the builder's bid commits to the tickets of slot $n$, keys are released once the payload is out and must be seen before the payload-timeliness committee votes, and the decrypted transactions execute at the top of block $n+1$, before any ordinary transaction. A short window shortens the delay the puzzle must provide and limits what a sender learns before deciding whether to reveal (\cref{sec:withholding}).

\subsection{The timing condition}\label{sec:timing}

The attacker's clock starts at \emph{exposure}, $\texp$, the first moment anyone but the sender holds the puzzle. The deadline $\torder$ is the moment the transaction's position can no longer change. The contents stay hidden until the deadline if
\[
  \torder - \texp < N \cdot \lattacker,
\]
with a margin for propagation. A forced opening finishes $N \cdot \lsolver$ after a solver starts, which on a laptop core is minutes after the deadline. If the attacker is $k$ times faster than assumed, the delay shrinks $k$-fold; once it falls below $\torder - \texp$, the contents can be read before the order is fixed. Privacy before ordering is then lost and content-dependent MEV becomes possible again, for whoever can solve fast enough.

\subsection{Choosing the delay}\label{sec:choosing-t}

The delay should be the smallest that covers the time from exposure to the fixed order, $\torder - \texp$. Building, forced opening, the solvers' load and the window in which a silent sender holds an option all grow linearly with $T$; only the size does not. The attacker speed already carries a margin of about 3.7 over the fastest published design (\cref{sec:pricing}), so we take $T$ about twice the time to fix the order. \Cref{tab:choosing-t} applies this to the settings of \cref{sec:realistic}. We use $T = 30$\,s as the reference: it covers a puzzle that is sent only to the slot's builder, or published after a commitment, and ordered within one 12-second slot.

\begin{table}[t]
  \centering\small
  \caption{The delay per setting. Building uses all cores of the M4 Air (0.70\,s per second of delay); forced opening uses one of its cores (6.6 times the delay).}
  \label{tab:choosing-t}
  \begin{tabularx}{\textwidth}{@{}>{\raggedright\arraybackslash}X>{\raggedright\arraybackslash}p{2.7cm}rrr@{}}
    \toprule
    Setting & Time to fix the order & Delay $T$ & Building & Forced opening \\
    \midrule
    Rollup with a sequencer & 1--2\,s & 15\,s & 11\,s & 1.7\,min \\
    Ethereum with enshrined PBS; puzzle sent to the builder, or published after a commitment & up to one slot, 12\,s & \textbf{30\,s} & 21\,s & 3.3\,min \\
    Sealed-bid auction; bid committed, then posted near the close & about one slot & 30\,s & 21\,s & 3.3\,min \\
    Public gossip, waiting $w$ slots & $12(w+1)$\,s & $24(w+1)$\,s & $17(w+1)$\,s & $2.7(w+1)$\,min \\
    Order fixed only at finality, Ethereum today & about 13\,min & 15\,min & 10.5\,min & 1.7\,h \\
    \bottomrule
  \end{tabularx}
\end{table}

The last row is why \leanseal needs an order that is fixed early. Waiting for today's finality makes every cost thirty times larger than in the reference setting, so a protocol either accepts that a reorganization exposes revealed transactions (\cref{sec:reorgs}) or relies on fast finality, which brings the delay back to the reference range. Public gossip is affordable for a wait of a slot or two but not more, which is why the puzzle should not wait in a public mempool.

\subsection{What \leanseal provides}\label{sec:mempool-props}

\Cref{tab:props} lists the properties an encrypted mempool may want and what \leanseal gives. No known design gives all of them; \cref{sec:comparison} compares the alternatives.

\begin{table}[t]
  \centering\small
  \caption{Properties of an encrypted mempool, and what \leanseal provides.}
  \label{tab:props}
  \begin{tabularx}{\textwidth}{@{}>{\raggedright\arraybackslash}p{4.3cm}>{\raggedright\arraybackslash}X@{}}
    \toprule
    Property & \leanseal \\
    \midrule
    Contents hidden until the order is fixed & yes, for $T$ after exposure, under the hardware assumption \\
    Every included transaction opens & yes, without the sender or anyone else, if the puzzle is well formed \\
    No free cancellation after seeing the order & no: withholding still delays or, with expiry, cancels execution; a deposit must deter it (\cref{sec:withholding}) \\
    Excluded transactions stay private & no: they open after $T$ like any other (\cref{sec:noninclusion}) \\
    Sender, fee and size hidden & no; the envelope is public and the size can be padded \\
    Trusted parties or setup & none \\
    Post-quantum & plausibly; proved for the idealized puzzle only (\cref{sec:attacks}) \\
    Cost & 152 to 512 bytes per transaction, sender work proportional to $T$ paid in advance, no work for the network in the normal case (\cref{sec:costs}) \\
    \bottomrule
  \end{tabularx}
\end{table}

\subsection{A proof of concept}\label{sec:poc}

We implemented this flow as a contract on any chain with EIP-7702, with off-chain sender, keeper and solver roles~\cite{leanSealRepo}. The block that includes a sealed transaction fixes its position within its batch. The sender reveals its key within two blocks, and anyone then executes the batch's revealed transactions in order: the contract checks each key against its commitment and decrypts the envelope, both with SHA-256 only, and calls the sender's account, which delegates to the contract for this. A missed reveal lets anyone who later finds the key from the puzzle open the transaction. On a local chain with 2-second blocks, a searcher watching the mempool sandwiched a plaintext swap, which lost 9.25\% against the fair price; the same swap, sealed, executed at the fair price. With its key withheld, a demonstration puzzle of $2^{29}$ steps, 1/16 of the reference work, was solved on one core in 19.5\,s and the swap executed late.

The prototype demonstrates submission, key revelation, ordered execution within a batch and forced disclosure. It does not demonstrate sandwich protection. Its searcher attacks only plaintext submissions. Once a key is revealed, anyone can decrypt the transaction and trade against the same pool before the separate execution call, and batches from different blocks can execute out of order. A deployment must reserve the execution position relative to every transaction that touches the same state, as LUCID does by executing decrypted transactions at the top of the next block; fixing the block that contains the submission is not enough. Nor does the prototype implement the deterrent: the deposit of a missed reveal goes to whoever posts the key, and the withholding sender knows it, so it can recover its own deposit. A deployment must separate an irrecoverable penalty from any solver reward. In the contract's unit tests, with a 128-byte test puzzle, submitting costs about 199{,}000 gas and revealing about 74{,}000, and a batch of two swaps executes for about 155{,}000 gas per swap including decryption. These are contract-call costs; calldata and intrinsic costs come on top and grow with the envelope.
